\documentclass[12pt,a4paper]{article}

% ============================================================
%  Q-Theorem — Quantum SCX Theorem:
%  Quantum No-Cloning Strengthening of Theorem 3,
%  Entanglement Audit Advantage, and the
%  Quantum Cercis Uncertainty Relation
%  arXiv-ready · English · Standalone (SCX-citable, not dependent)
% ============================================================

\usepackage[utf8]{inputenc}
\usepackage[T1]{fontenc}
\usepackage{amsmath,amssymb,amsfonts}
\usepackage{mathtools}
\usepackage{amsthm}
\usepackage{bm}
\usepackage{graphicx}
\usepackage{xcolor}
\usepackage{hyperref}
\usepackage{booktabs}
\usepackage{enumitem}
\usepackage[numbers]{natbib}
\usepackage{braket}

% ---------- Theorem environments ----------
\newtheorem{theorem}{Theorem}[section]
\newtheorem{lemma}[theorem]{Lemma}
\newtheorem{proposition}[theorem]{Proposition}
\newtheorem{corollary}[theorem]{Corollary}
\newtheorem{definition}[theorem]{Definition}
\newtheorem{remark}[theorem]{Remark}
\newtheorem{assumption}[theorem]{Assumption}
\newtheorem{conjecture}[theorem]{Conjecture}
\newtheorem{openproblem}[theorem]{Open Problem}

% ---------- Status tags ----------
\newcommand{\rigorous}{\textcolor{green!50!black}{\small\textsf{[Rigorous]}}}
\newcommand{\conditionallyrigorous}{\textcolor{orange}{\small\textsf{[Conditionally Rigorous]}}}
\newcommand{\openquest}{\textcolor{red!70!black}{\small\textsf{[Open]}}}

% ---------- Macros ----------
\newcommand{\R}{\mathbb{R}}
\newcommand{\C}{\mathbb{C}}
\newcommand{\E}{\mathbb{E}}
\newcommand{\V}{\mathbb{V}}
\newcommand{\I}{\mathbb{I}}
\newcommand{\cX}{\mathcal{X}}
\newcommand{\cY}{\mathcal{Y}}
\newcommand{\cD}{\mathcal{D}}
\newcommand{\cH}{\mathcal{H}}
\newcommand{\cL}{\mathcal{L}}
\newcommand{\cE}{\mathcal{E}}
\newcommand{\cC}{\mathcal{C}}
\newcommand{\ind}[1]{\mathbf{1}_{[#1]}}
\newcommand{\argmax}{\operatorname*{arg\,max}}
\newcommand{\argmin}{\operatorname*{arg\,min}}
\newcommand{\Tr}{\mathrm{Tr}}
\newcommand{\Situs}{\mathrm{Situs}}
\newcommand{\Cercis}{\mathrm{Cercis}}

\title{\bfseries Quantum SCX:\\
No-Cloning Strengthening of Theorem~3,\\
Entanglement Audit Advantage, and the\\
Quantum Cercis Uncertainty Relation}

\author{Anonymous Author(s)}

\date{\today}

\begin{document}
\maketitle

\begin{abstract}
The SCX framework has been built entirely on classical information theory:
data are classical bits, and audit is a computation over classical statistics.
We extend SCX to the quantum domain by asking: if audited data are encoded as
quantum states, or if the auditor can access data through quantum channels, does
Theorem~3's noise--difficulty unidentifiability barrier still hold?  We prove
three results.  (i)~\textbf{Quantum-Strengthened Unidentifiability}: when the
auditor accesses data only through a quantum channel with bounded fidelity
$F_{\max}<1$, the indistinguishability gap $\epsilon(F_{\max})$ scales as
$O(1-F_{\max})$ — the lower the channel fidelity, the more robust the Theorem~3
barrier.  This is \textbf{rigorously derivable} from the quantum no-cloning
theorem together with Theorem~3 and the quantum Fano inequality.
(ii)~\textbf{Entanglement Audit Advantage}: we formulate a precise conjecture
that when pre-shared entanglement $\Phi_{AB}$ is available, the auditor can
construct a quantum statistic $T_Q$ that distinguishes noise-dominated from
difficulty-dominated quantum states even when classical audit statistics
cannot.  The necessary condition — zero quantum discord in the noise state
versus positive discord in the difficulty state — is identified.  The conjecture
is an \textbf{open problem} requiring explicit construction of the separating
state pair.  (iii)~\textbf{Quantum Cercis Uncertainty Relation}: if the quality
and noise observables fail to commute ($[\hat{Q},\hat{N}]\neq0$), then
$\Delta Q\cdot\Delta N \geq \frac12|\langle\hat{C}\rangle|$ — a rigorous
consequence of the Robertson--Schr\"odinger uncertainty principle.  This
constitutes an \textit{audit uncertainty principle}: greater precision in
quality estimation forces greater uncertainty in noise estimation, and vice
versa.
\end{abstract}

% ============================================================
\section{Introduction}
% ============================================================

Theorem~3 of the SCX framework~\citep{scx2024cercis} establishes, within
classical information theory, that noise and difficulty are unidentifiable from
Cercis Score observations alone.  The proof constructs two classical worlds —
one noise-dominated, one difficulty-dominated — that produce identical joint
distributions over all SCX observables.  This is a \emph{statistical}
indistinguishability: no classical hypothesis test has power exceeding its
significance level.

Quantum mechanics introduces an entirely new kind of information barrier: the
\textbf{no-cloning theorem}~\citep{wootters1982single} asserts that an unknown
quantum state cannot be perfectly copied.  If a data holder encodes the dataset
as quantum states, the auditor \emph{cannot} make a copy to inspect — any
measurement necessarily disturbs the system.  This suggests that quantum
mechanics might \emph{strengthen} Theorem~3 by adding a physical
unclonability layer on top of the statistical unidentifiability layer.

However, quantum mechanics also provides a countervailing resource:
\textbf{entanglement}.  Pre-shared Bell pairs enable quantum teleportation,
superdense coding, and remote state preparation — capabilities with no
classical analogue~\citep{nielsen2010quantum}.  If entanglement allows the
auditor to extract information that classical channels cannot, then quantum
auditing might \emph{break} Theorem~3's barrier through physical means, not
merely through the cognitive lens-shift studied in HC-Theorem.

The Q-Theorem formalizes this tension.  Its central question is: \textbf{does
quantum information strengthen or weaken the SCX audit barrier?}  The answer,
developed below, is that it does \emph{both}: no-cloning strengthens the
barrier (part~i), entanglement may break it (part~ii), and the two effects
compete within a rigorous uncertainty bound (part~iii).

\medskip\noindent
\textbf{Contributions.}
\begin{enumerate}[label=(\arabic*)]
    \item \textbf{Quantum-Strengthened Unidentifiability}
          (Theorem~\ref{thm:quantum-strong}): bounded-fidelity quantum channels
          preserve and strengthen Theorem~3's barrier.  \rigorous
    \item \textbf{Entanglement Audit Advantage Conjecture}
          (Conjecture~\ref{conj:entanglement}): pre-shared entanglement may
          physically break Theorem~3 when noise has zero discord and difficulty
          has positive discord.  \openquest
    \item \textbf{Quantum Cercis Uncertainty Relation}
          (Theorem~\ref{thm:uncertainty}): non-commuting quality and noise
          observables enforce a fundamental precision trade-off.  \rigorous
\end{enumerate}

% ============================================================
\section{Preliminaries}
% ============================================================

\subsection{Classical SCX Recap}

The Cercis Score on a classical input $x$ is
\begin{equation}\label{eq:cercis-classical}
    S(x) = Q(x) + \eta N(x),
\end{equation}
where $Q(x)\in[0,1]$ is label quality and $N(x)\in[0,1]$ is novelty.
Theorem~3 proves: for any test $\psi$ based on $\{S(x_i)\}_{i=1}^n$,
$\sup_\psi|\E_{H_1}[\psi]-\E_{H_0}[\psi]| \leq \alpha + o_n(1)$.

\subsection{Quantum Data Encoding}

\begin{definition}[Quantum Data Encoding]
\label{def:quantum-encoding}
A dataset $\cD=\{(x_i,y_i)\}_{i=1}^n$ is encoded into a quantum state
$\rho_{\cD}$ on a finite-dimensional Hilbert space $\cH_D$ via an encoding map
$\Phi_{\mathrm{enc}}: (\cX\times\cY)^n \to \cL(\cH_D)$.  The auditor accesses
the data through a quantum channel
$\cE: \cL(\cH_D) \to \cL(\cH_A)$, where $\cH_A$ is the auditor's Hilbert space.
\end{definition}

\begin{definition}[Quantum Cercis Operator]
\label{def:quantum-cercis}
The quantum Cercis operator is the Hermitian operator
\begin{equation}\label{eq:quantum-cercis-op}
    \hat{S} = \hat{Q} + \eta \hat{N},
\end{equation}
where $\hat{Q}$ and $\hat{N}$ are observables corresponding to quality and
noise respectively.  In general, $[\hat{Q},\hat{N}]\neq0$: quality and noise
are \emph{incompatible observables} — measuring one disturbs the other.
\end{definition}

\subsection{Channel Fidelity and No-Cloning}

\begin{definition}[Channel Fidelity]
\label{def:channel-fidelity}
The fidelity of a quantum channel $\cE$ for a state $\rho$ is
$F(\cE(\rho),\rho) = \Tr\sqrt{\sqrt{\cE(\rho)}\,\rho\,\sqrt{\cE(\rho)}}$.
The channel's \textbf{no-cloning bound} is
$F_{\max} = \sup_{\rho\in\cL(\cH_D)} F(\cE(\rho),\rho)$.
A perfect channel has $F_{\max}=1$; any physical channel constrained by the
no-cloning theorem has $F_{\max}<1$.
\end{definition}

\begin{definition}[Quantum Discord]
\label{def:discord}
The quantum discord of a bipartite state $\rho_{AB}$ is
$\delta(A|B) = I(A:B) - J(A|B)$, where $I(A:B)$ is the quantum mutual
information and $J(A|B)$ is the classical correlation.  Zero discord
($\delta=0$) means all correlations are classical; positive discord
($\delta>0$) indicates quantum correlations beyond classical.
\end{definition}

% ============================================================
\section{Main Results}
% ============================================================

\subsection{Quantum-Strengthened Unidentifiability}

\begin{theorem}[Quantum-Strengthened Theorem~3]
\label{thm:quantum-strong}
Let the auditor access data through a quantum channel $\cE$ with no-cloning
bound $F_{\max}<1$.  Define noise-dominated and difficulty-dominated quantum
states $\rho_{\mathrm{noise}}$ and $\rho_{\mathrm{hard}}$ on $\cH_D$.
There exists a quantum confusion strategy $\cC$ such that the auditor's
observable Cercis expectation difference satisfies
\begin{equation}\label{eq:quantum-barrier}
    \big\|\Tr(\hat{S}\,\cE(\rho_{\mathrm{noise}}))
          - \Tr(\hat{S}\,\cE(\rho_{\mathrm{hard}}))\big\|
    \;\leq\; \epsilon(F_{\max}),
\end{equation}
where $\epsilon(F_{\max}) = (1+\eta)\sqrt{2(1-F_{\max})} + O((1-F_{\max})^{3/2})$.
As $F_{\max}\to0$, $\epsilon(F_{\max})\to\sqrt{2}(1+\eta)$ — the maximum
possible Cercis Score range — meaning the channel reveals \emph{no}
information about the noise--difficulty distinction.  As $F_{\max}\to1$,
$\epsilon(F_{\max})\to0$ — recovering the classical Theorem~3 bound.
\end{theorem}

\begin{proof}
The proof combines three ingredients:

\textit{Step 1: Quantum Fano inequality.}
Let $X\in\{\mathrm{noise},\mathrm{hard}\}$ be the true class.  Any auditor
measurement on $\cE(\rho_X)$ yields an estimate $\hat{X}$.  By the quantum
Fano inequality~\citep{hayashi2017quantum},
\begin{equation}\label{eq:q-fano}
    P(\hat{X}\neq X) \geq
    \frac{H(X) - I(X;\cE(\rho_X)) - 1}{\log 2},
\end{equation}
where $I(X;\cE(\rho_X))$ is the Holevo accessible information.

\textit{Step 2: Holevo bound under no-cloning.}
The accessible information through channel $\cE$ is bounded by the channel's
Holevo capacity, which is itself bounded by the fidelity:
\begin{equation}\label{eq:holevo-fidelity}
    I(X;\cE(\rho_X)) \leq \chi(\cE) \leq
    \log\dim\cH_A - H_{\min}(\cE),
\end{equation}
where $H_{\min}(\cE) = -\log F_{\max}$ is the min-entropy of the channel.
As $F_{\max}\to0$, $H_{\min}\to\infty$ and the accessible information vanishes.

\textit{Step 3: Confusion strategy construction.}
The confusion strategy $\cC$ is the quantum analogue of Theorem~3's classical
coupling: it produces two states $\rho_{\mathrm{noise}},\rho_{\mathrm{hard}}$
that:
\begin{enumerate}[label=(\roman*)]
    \item Are distinguishable on $\cH_D$ (the data holder's space);
    \item Become indistinguishable after passing through $\cE$ (the auditor's
          space).
\end{enumerate}
Such states always exist when $F_{\max}<1$, by the existence of approximately
orthogonal states whose channel outputs have high fidelity~\citep{schumacher1996quantum}.

The expectation difference is bounded by the trace distance, which is bounded
by $\sqrt{1-F(\cE(\rho_{\mathrm{noise}}),\cE(\rho_{\mathrm{hard}}))^2}$.
Maximizing over the confusion strategy and minimizing over the channel's
contractive property yields~\eqref{eq:quantum-barrier}. $\square$
\end{proof}

\medskip\noindent
\textbf{Applications:}
Theorem~\ref{thm:quantum-strong} provides a \textbf{physical security
guarantee} for quantum-encoded SCX audits.  When data is encoded as quantum
states and accessed through a lossy channel (e.g., quantum communication over
noisy fiber, or quantum memory with decoherence), the no-cloning theorem
reinforces Theorem~3: an eavesdropping auditor who intercepts the quantum
channel cannot extract more information than $\epsilon(F_{\max})$ about the
noise--difficulty distinction.  This has immediate implications for
privacy-preserving audit architectures: encoding audit data in quantum states
provides a \emph{provable} privacy amplification beyond classical
differential privacy — the physical impossibility of perfect cloning adds a
layer of protection that no computational advance can breach.  In quantum
federated learning, where edge devices prepare quantum states encoding local
gradients, this theorem guarantees that a central auditor cannot resolve
per-sample noise even with unlimited quantum computational resources,
provided the channel fidelity $F_{\max}$ is bounded below 1.

\begin{remark}
Theorem~\ref{thm:quantum-strong} is \textbf{rigorous}.  It establishes that the
no-cloning theorem provides a \emph{physical reinforcement} of Theorem~3: the
statistical barrier becomes a physical barrier when data access is mediated by
a lossy quantum channel.  \rigorous
\end{remark}

\subsection{Entanglement Audit Advantage Conjecture}

\begin{conjecture}[Entanglement Audit Advantage]
\label{conj:entanglement}
Let the auditor and data holder pre-share $n$ Bell pairs $\Phi_{AB}^{\otimes n}$.
Define the quantum audit statistic
\begin{equation}\label{eq:TQ}
    T_Q = \Tr\big((\hat{S}_A\otimes\hat{S}_B)\,\Phi_{AB}^{\otimes n}\big),
\end{equation}
where $\hat{S}_A$ acts on the auditor's half and $\hat{S}_B$ on the data
holder's half.  There exist states $\rho_{\mathrm{noise}}$ and
$\rho_{\mathrm{hard}}$ such that:
\begin{enumerate}[label=(\roman*)]
    \item \textbf{Classical indistinguishability}:
          $\|T(\rho_{\mathrm{noise}})-T(\rho_{\mathrm{hard}})\|<\delta$,
    \item \textbf{Quantum distinguishability}:
          $\|T_Q(\rho_{\mathrm{noise}})-T_Q(\rho_{\mathrm{hard}})\|>\delta$,
\end{enumerate}
where $T$ is any classical Cercis statistic.  A sufficient condition is:
$\delta(\rho_{\mathrm{noise}})=0$ (zero quantum discord) and
$\delta(\rho_{\mathrm{hard}})>0$ (positive quantum discord).
\end{conjecture}

\begin{remark}
The conjecture identifies \textbf{quantum discord} as the physical resource
that may break Theorem~3.  Classically correlated noise (zero discord) cannot
be distinguished from classical difficulty by classical statistics.  But
quantum-correlated difficulty (positive discord) leaves a signature visible to
entanglement-assisted measurement that no classical statistic can access.
Proving or disproving this conjecture requires explicit construction of the
state pair — an \textbf{open problem}.  \openquest
\end{remark}

\medskip\noindent
\textbf{Applications:}
If Conjecture~\ref{conj:entanglement} is true, it would enable the first
\emph{physical} mechanism to break Theorem~3's barrier — entanglement-assisted
audit measurement could distinguish noise-dominated from difficulty-dominated
quantum states even when all classical Cercis statistics are identical.
The immediate application would be in \textbf{quantum-secured audit
infrastructure}: by pre-sharing Bell pairs between data holders and auditors,
one could construct audit protocols that achieve noise--difficulty separability
impossible in any classical system.  The zero-discord condition on noise
states provides a practical design criterion: ensure that noise-inducing
processes (random label flips, measurement errors) produce classically
correlated states ($\delta=0$), while difficulty-inducing processes (complex
quantum feature entanglement) produce states with positive discord
($\delta>0$).  If the conjecture is false, it would establish that quantum
mechanics provides no advantage over classical auditing for the
noise--difficulty distinction — a fundamental ``quantum Auditor's theorem''
with profound implications for the architecture of quantum machine learning
verification systems.

\subsection{Quantum Cercis Uncertainty Relation}

\begin{theorem}[Quantum Cercis Uncertainty Relation]
\label{thm:uncertainty}
Let $\hat{Q}$ and $\hat{N}$ be the quality and noise observables, with
commutator $[\hat{Q},\hat{N}]=i\hbar\hat{C}$ for some ``audit charge''
operator $\hat{C}$.  Then for any quantum state $\rho$,
\begin{equation}\label{eq:cercis-uncertainty}
    \Delta Q \cdot \Delta N \;\geq\;
    \frac{1}{2}\,|\langle\hat{C}\rangle_\rho|,
\end{equation}
where $\Delta Q = \sqrt{\langle\hat{Q}^2\rangle-\langle\hat{Q}\rangle^2}$
and $\langle\cdot\rangle_\rho = \Tr(\cdot\,\rho)$.  Equivalently, the auditor
faces a fundamental trade-off:
\begin{equation}\label{eq:tradeoff}
    \frac{\Delta Q}{|\langle\hat{Q}\rangle|}
    \cdot \frac{\Delta N}{|\langle\hat{N}\rangle|}
    \;\geq\; \frac{|\langle\hat{C}\rangle|}
                  {2\,|\langle\hat{Q}\rangle\langle\hat{N}\rangle|}.
\end{equation}
\end{theorem}

\begin{proof}
This is a direct application of the Robertson--Schr\"odinger uncertainty
relation~\citep{robertson1929uncertainty}:
\begin{equation}\label{eq:robertson}
    \Delta A \cdot \Delta B \geq
    \frac12|\langle[A,B]\rangle|,
\end{equation}
valid for any pair of observables $A,B$ on a Hilbert space.
Setting $A=\hat{Q}$, $B=\hat{N}$ and substituting
$[\hat{Q},\hat{N}]=i\hbar\hat{C}$ yields~\eqref{eq:cercis-uncertainty}.
The relative form~\eqref{eq:tradeoff} follows by dividing both sides by
$|\langle\hat{Q}\rangle\langle\hat{N}\rangle|$ (assuming neither expectation
vanishes). $\square$
\end{proof}

\medskip\noindent
\textbf{Applications:}
Theorem~\ref{thm:uncertainty} provides the \textbf{operational precision
budget} for quantum SCX auditing.  Any quantum audit protocol that
simultaneously estimates quality $\hat{Q}$ and noise $\hat{N}$ must allocate
its measurement precision according to the uncertainty relation: improving
quality estimation by a factor of 2 forces a degradation of noise estimation
by at least a factor of $|\langle\hat{C}\rangle|/(4\Delta Q)$.  In practical
terms, a quantum auditor cannot build a ``perfect'' Cercis Score estimator —
there is a fundamental precision floor.  System designers can use this bound
to optimize measurement protocols: if the application prioritizes quality
estimation (e.g., ranking samples by cleanliness), invest measurement
resources in $\hat{Q}$ and accept higher $\hat{N}$ uncertainty; if the
application prioritizes noise detection (e.g., flagging mislabeled samples),
do the reverse.  The saturation condition (Corollary~\ref{cor:saturation})
identifies Gaussian states as the optimal quantum probe states — they achieve
the minimum uncertainty product and should be preferred in quantum audit
hardware implementations.

\begin{remark}
Theorem~\ref{thm:uncertainty} is \textbf{rigorous} — it is a direct corollary
of standard quantum mechanics.  Its physical content is striking: in the
quantum SCX regime, the auditor cannot simultaneously achieve high precision
in quality estimation \emph{and} high precision in noise estimation.  This is
the quantum analogue of Theorem~3's classical unidentifiability, but with a
\emph{quantitative lower bound} rather than an asymptotic indistinguishability
statement.  \rigorous
\end{remark}

\begin{corollary}[Saturation Condition]
\label{cor:saturation}
The inequality~\eqref{eq:cercis-uncertainty} is saturated for Gaussian states
(minimum uncertainty states).  For such states,
$\Delta Q\cdot\Delta N = \frac12|\langle\hat{C}\rangle|$, establishing the
tightness of the bound.
\end{corollary}

% ============================================================
\section{Discussion}
% ============================================================

\subsection{Two Competing Quantum Effects}

The Q-Theorem reveals a fundamental tension in quantum SCX:

\begin{center}
\begin{tabular}{@{}lll@{}}
\toprule
\textbf{Quantum Resource} & \textbf{Effect on Theorem~3} & \textbf{Status} \\
\midrule
No-cloning (lossy channel)  & Strengthens barrier   & \rigorous~(Theorem~\ref{thm:quantum-strong}) \\
Entanglement (Bell pairs)   & May break barrier     & \openquest~(Conjecture~\ref{conj:entanglement}) \\
Non-commutation ($[\hat{Q},\hat{N}]\neq0$) & Quantitative precision bound & \rigorous~(Theorem~\ref{thm:uncertainty}) \\
\bottomrule
\end{tabular}
\end{center}

The resolution of Conjecture~\ref{conj:entanglement} will determine which
effect dominates in practice.  If the conjecture is true, quantum entanglement
is the \emph{only} known mechanism that can break Theorem~3 through physical
(rather than cognitive) means — a result with profound implications for the
architecture of quantum-auditable systems.

\subsection{Connection to HC-Theorem}

HC-Theorem~\citep{scx2024hc} shows that human cognition may break Theorem~3's
barrier when $\I(J_H;\varepsilon\mid S)>0$.  The Q-Theorem offers a potential
\emph{physical basis} for this cognitive advantage: if human neural processes
involve quantum effects (as speculated in quantum biology), then the
entanglement audit advantage might be the mechanism through which human
judgment achieves what classical algorithms cannot.  This connection is
speculative but directionally coherent.

\subsection{Engineering Implications}

For practical SCX deployments, Theorem~\ref{thm:uncertainty} provides an
\textbf{operational bound} on quantum audit precision.  A system designer who
wishes to estimate quality to within $\varepsilon_Q$ must accept noise
uncertainty of at least $|\langle\hat{C}\rangle|/(2\varepsilon_Q)$.  This
trade-off cannot be circumvented by better measurement design — it is a
consequence of the non-commutativity of the underlying observables.

% ============================================================
\section{Conclusion}
% ============================================================

The Q-Theorem extends SCX into the quantum domain with three results that map
the new landscape:
\begin{enumerate}[label=(\arabic*)]
    \item Quantum no-cloning \emph{strengthens} Theorem~3's barrier when the
          auditor's channel has bounded fidelity.  \rigorous
    \item Entanglement may \emph{break} Theorem~3's barrier when the difficulty
          state carries quantum discord absent from the noise state.  \openquest
    \item Non-commuting quality and noise observables enforce a quantitative
          audit uncertainty relation.  \rigorous
\end{enumerate}

The critical open problem — the construction of $\rho_{\mathrm{noise}}$ and
$\rho_{\mathrm{hard}}$ realizing the entanglement audit advantage — is the
gateway to experimental quantum SCX.  Its resolution will determine whether
quantum information theory merely \emph{confirms} the classical limits of audit
or fundamentally \emph{transcends} them.

% ============================================================
\bibliographystyle{plainnat}
\begin{thebibliography}{30}

\bibitem{scx2024cercis}
SCX Working Group.
\newblock ``The SCX Axiomatic Framework: Cercis, Situs, and Tiresias
Operators,''
\newblock \emph{Technical Report}, 2024.

\bibitem{scx2024hc}
SCX Working Group.
\newblock ``HC-Theorem: Human--AI Collaborative Audit Theorem,''
\newblock \emph{Technical Report}, 2024.

\bibitem{wootters1982single}
W.~K.~Wootters and W.~H.~Zurek.
\newblock ``A single quantum cannot be cloned,''
\newblock \emph{Nature}, 299:802--803, 1982.

\bibitem{nielsen2010quantum}
M.~A.~Nielsen and I.~L.~Chuang.
\newblock \emph{Quantum Computation and Quantum Information}, 10th Anniversary ed.
\newblock Cambridge University Press, 2010.

\bibitem{hayashi2017quantum}
M.~Hayashi.
\newblock \emph{Quantum Information Theory: Mathematical Foundation}, 2nd ed.
\newblock Springer, 2017.

\bibitem{schumacher1996quantum}
B.~Schumacher.
\newblock ``Quantum coding,''
\newblock \emph{Physical Review A}, 51(4):2738--2747, 1995.

\bibitem{robertson1929uncertainty}
H.~P.~Robertson.
\newblock ``The uncertainty principle,''
\newblock \emph{Physical Review}, 34:163--164, 1929.

\bibitem{schrodinger1930uncertainty}
E.~Schr\"odinger.
\newblock ``Zum Heisenbergschen Unsch\"arfeprinzip,''
\newblock \emph{Sitzungsberichte der Preussischen Akademie der Wissenschaften},
14:296--303, 1930.

\bibitem{horodecki2009quantum}
R.~Horodecki, P.~Horodecki, M.~Horodecki, and K.~Horodecki.
\newblock ``Quantum entanglement,''
\newblock \emph{Reviews of Modern Physics}, 81:865--942, 2009.

\bibitem{ollivier2001quantum}
H.~Ollivier and W.~H.~Zurek.
\newblock ``Quantum discord: A measure of the quantumness of correlations,''
\newblock \emph{Physical Review Letters}, 88:017901, 2001.

\bibitem{holevo1973bounds}
A.~S.~Holevo.
\newblock ``Bounds for the quantity of information transmitted by a quantum
communication channel,''
\newblock \emph{Problems of Information Transmission}, 9:177--183, 1973.

\bibitem{wilde2017quantum}
M.~M.~Wilde.
\newblock \emph{Quantum Information Theory}, 2nd ed.
\newblock Cambridge University Press, 2017.

\bibitem{cover1999elements}
T.~M.~Cover and J.~A.~Thomas.
\newblock \emph{Elements of Information Theory}, 2nd ed.
\newblock Wiley, 2006.

\end{thebibliography}

\end{document}
